% 05-discussion — v4.0
\subsection*{What survived, and what did not}
The two halves of the benchmark give different answers. Alpha reactivity survives the
acquisition chain almost intact: at the same electrode, the chain costs 0.01\% of the
effect. What costs alpha is the choice of electrode. Fz has about 58\% of the
eyes-closed alpha power of a posterior average, and this difference is a property of the
scalp, not of the electronics. Noise also matters, and it matters in a predictable way:
the effect declines steadily and reaches 59\% of the reference at 0\,dB. Bandwidth,
sampling rate and quantisation, the specifications usually printed on a consumer
headset's box, changed the 8--12\,Hz estimate by a negligible amount.

Affective decoding behaves differently. With the full montage, valence, arousal and
four-quadrant classification are above chance. With one electrode they are not: valence
is undetectable, arousal is weak, and four-class accuracy falls to 0.34. The contrast is
not a statement about broad alpha versus distributed affect. The two analyses use
different effect sizes, different numbers of trials and different estimators, and the
alpha ratio is partly invariant to the operations we applied. The useful conclusion is
narrower: within this benchmark, the acquisition chain does not limit alpha, but it does
limit affective decoding once the spatial representation is reduced to one site.

\subsection*{Where the affective result goes}
The decomposition helps. Resampling, quantisation and band-limiting leave full-montage
accuracy unchanged. The loss appears when the data become one electrode, and it appears
at all four midline sites we tried. Valence is not distinguishable from chance at Fz,
Cz, Pz or Oz, and four-class accuracy stays between 0.33 and 0.38. Arousal is the
exception: it remains weakly above chance at Fz, Pz and Oz before correction, though not
after correction across the site sweep. We treat the site sweep as exploratory.

The band control changes how the montage result should be read. With only theta and alpha
features, valence falls to 0.60 and four-class accuracy to 0.46. Beta-only features give
0.70 and 0.62, and beta-only is at least as good as the full montage. Beta is also where
EMG contamination is most likely \cite{whitham2007emg,goncharova2003emg}. We did not
record EMG or test a paired beta-versus-theta/alpha difference, so the result is an
interpretation: part of the montage advantage may be muscular rather than cortical. That
possibility should be ruled out before the montage number is used as a neural benchmark.

The low-cost Emotiv arm is a useful sanity check. Its accuracy is close to the high-end
montage on the 12 participants present in both datasets, and the matched differences are
small. The equivalence test does not pass, however, so the appropriate description is
"similar in this sample", not "equivalent". The comparison also mixes channel count,
feature count, recording units, epoch window and condition arm, and eight participants
are missing from the Emotiv analysis. It is a check, not a device validation.

\subsection*{What this means for reporting}
The first lesson is about reference choice. Comparing one electrode with a whole-head
average confounds channel count with electrode site, and a benchmark that does so will
understate what a single channel retains. The correction is simple: compare the same
site before and after degradation. Reference and site dependence of scalp power is not a
new observation \cite{hagemann2001reference}; what is new here is the size of the effect
for a single-channel benchmark, and the recommendation to make same-site comparison the
default.

The second lesson is about the site itself. Fz is a convenient research midline
electrode, but most single-electrode consumer devices record at the forehead: Fp1 in
NeuroSky/MindWave and AF7/AF8 in Muse \cite{krigolson2017muse,cannard2021muse}. Frontal
sites have less posterior alpha and more ocular and muscle contamination, so the 42\%
site cost we measured is probably an underestimate for those devices. The consumer-EEG
reliability comparison we cite is also at Fp1 \cite{rogers2016reliability}, not Fz. The
single-channel affective failure is established here for Fz, Cz, Pz and Oz, not for every
possible electrode.

The third lesson is about timescale. A 2\,s alpha window is a noisy estimate of a
momentary state; the odd--even correlation is a short-lag stability measure, not a
reliability coefficient, and its value near zero should not be read as "alpha is
unreliable". Block-level alpha power was much more stable (ICC 0.92). A product that
reports "your alpha right now" and a study that reports a stable individual difference
are making different claims.

\subsection*{Limitations}
The alpha arm uses all 109 subjects, but only the first 30 have the repeated motor runs
needed for the ICC estimate; the ICC therefore rests on a small subset and has no
confidence interval. The affective arm uses 18 high-end recordings, and the Emotiv arm
covers 12 participants. The Emotiv signal is in device units, its epoch window differs
slightly from the BrainVision epochs, and the archive's two condition arms are not
distinguished. We did not measure a real single-channel device, so the single-channel
condition is a simulation of that device, not the device itself. The blink model is
additive, spectrally orthogonal to alpha, and applied equally to both conditions; it
cannot test the real ocular interaction, and the pre-specified blink hypothesis is
therefore untested. Artifact handling is minimal: no ICA or Autoreject, so the beta-band
features may be contaminated. The affective labels are normative IAPS categories rather
than participants' own ratings, and we did not test picture-level generalisation or a
label-permutation null. Only four midline sites were swept, and the single-channel
condition has three features against 198 for the full montage, so the representation
comparison confounds information with dimensionality. The 0\,dB level is defined on
broadband variance, the in-band signal-to-noise ratio is not reported, and the ratio
estimator can compress the effect. Eight of 109 subjects had a negative reference
reactivity, which is expected for a resting-state contrast but limits the precision of
the mean. The analysis plan was not registered; several pre-specified items (Fp1/Cz
sensitivity, the 60\% floor, AUC and median-split outcomes) were changed or dropped, and
the study is exploratory in the strict sense. The motor arm is invalid, not merely
underpowered, because its reference decoder did not exceed chance. Finally, the
equivalence test uses a small sample and a margin that is wide relative to the observed
differences.

\subsection*{Where to go next}
The obvious next step is a real single-channel device. A small study with a consumer
headset, raw samples and the same pipeline would show whether the simulation is
realistic, and whether the site cost is larger at the forehead, as we expect. A second
step is to measure EMG or add a high-frequency control to separate muscular from
cortical contributions to the beta-band result. A third is to register the analysis
plan and repeat the alpha arm with a pre-specified decision rule. None of these requires
new theory; they require better measurement, which is the point of the benchmark.
